\documentclass[10pt,a4paper]{amsart}
\usepackage{amsmath,amssymb,amsthm,graphicx}
% Retained minimal preamble for the Stacks Project chapter sites-cohomology.tex (Cohomology on Sites).
% Verbatim excerpt of the chapter preamble: Xy-pic setup (native preamble lines 14-18) and the native
% theorem declarations with the mathematical macros (native preamble lines 154-217). Copied mechanically;
% nothing added, removed or redefined. The wrapper class and the amsmath/amssymb/amsthm/graphicx base
% packages are supplied by the assembler, so the class/externaldocument part of the native preamble is omitted.
% For commutative diagrams we use Xy-pic
\usepackage[all]{xy}

% We use 2cell for 2-commutative diagrams.
\xyoption{2cell}
% Theorem environments.
%
\theoremstyle{plain}
\newtheorem{theorem}[subsection]{Theorem}
\newtheorem{proposition}[subsection]{Proposition}
\newtheorem{lemma}[subsection]{Lemma}

\theoremstyle{definition}
\newtheorem{definition}[subsection]{Definition}
\newtheorem{example}[subsection]{Example}
\newtheorem{exercise}[subsection]{Exercise}
\newtheorem{situation}[subsection]{Situation}

\theoremstyle{remark}
\newtheorem{remark}[subsection]{Remark}
\newtheorem{remarks}[subsection]{Remarks}

\numberwithin{equation}{subsection}

% Macros
%
\def\lim{\mathop{\mathrm{lim}}\nolimits}
\def\colim{\mathop{\mathrm{colim}}\nolimits}
\def\Spec{\mathop{\mathrm{Spec}}}
\def\Hom{\mathop{\mathrm{Hom}}\nolimits}
\def\Ext{\mathop{\mathrm{Ext}}\nolimits}
\def\SheafHom{\mathop{\mathcal{H}\!\mathit{om}}\nolimits}
\def\SheafExt{\mathop{\mathcal{E}\!\mathit{xt}}\nolimits}
\def\Sch{\mathit{Sch}}
\def\Mor{\mathop{\mathrm{Mor}}\nolimits}
\def\Ob{\mathop{\mathrm{Ob}}\nolimits}
\def\Sh{\mathop{\mathit{Sh}}\nolimits}
\def\NL{\mathop{N\!L}\nolimits}
\def\CH{\mathop{\mathrm{CH}}\nolimits}
\def\proetale{{pro\text{-}\acute{e}tale}}
\def\etale{{\acute{e}tale}}
\def\QCoh{\mathit{QCoh}}
\def\Ker{\mathop{\mathrm{Ker}}}
\def\Im{\mathop{\mathrm{Im}}}
\def\Coker{\mathop{\mathrm{Coker}}}
\def\Coim{\mathop{\mathrm{Coim}}}

% Boxtimes
%
\DeclareMathSymbol{\boxtimes}{\mathbin}{AMSa}{"02}

%
% Macros for moduli stacks/spaces
%
\def\QCohstack{\mathcal{QC}\!\mathit{oh}}
\def\Cohstack{\mathcal{C}\!\mathit{oh}}
\def\Spacesstack{\mathcal{S}\!\mathit{paces}}
\def\Quotfunctor{\mathrm{Quot}}
\def\Hilbfunctor{\mathrm{Hilb}}
\def\Curvesstack{\mathcal{C}\!\mathit{urves}}
\def\Polarizedstack{\mathcal{P}\!\mathit{olarized}}
\def\Complexesstack{\mathcal{C}\!\mathit{omplexes}}
% \Pic is the operator that assigns to X its picard group, usage \Pic(X)
% \Picardstack_{X/B} denotes the Picard stack of X over B
% \Picardfunctor_{X/B} denotes the Picard functor of X over B
\def\Pic{\mathop{\mathrm{Pic}}\nolimits}
\def\Picardstack{\mathcal{P}\!\mathit{ic}}
\def\Picardfunctor{\mathrm{Pic}}
\def\Deformationcategory{\mathcal{D}\!\mathit{ef}}
\usepackage{cleveref}
\crefname{originallemma}{original lemma}{original lemmas}
\Crefname{originallemma}{Original lemma}{Original lemmas}
\makeatletter
\newlabel{sites-modules-lemma-adjoint-hom-restrict}{{0932}{source}{}{sourcepointer-1}{}}
\newlabel{sites-modules-lemma-adjoint-hom-restrict@cref}{{[originallemma][0][]0932}{[1][0][]source}{}{}{}}
\makeatother
\title{Restriction of scalars along a map of sheaves of rings preserves K-injectivity}
\author{The Stacks Project Authors}
\date{stacks-project:sites-cohomology.tex:local-snapshot-20261011; GNU FDL 1.2 or later}
\begin{document}
\maketitle
\noindent\textit{Editorial note.} All delivered mathematical text is the author's own, copied from the pinned Stacks Project chapter file sites-cohomology.tex; no statement or proof step was rewritten, reordered or removed.\par\medskip
\noindent\textit{Editorial note.} The proof is the author's complete argument: it displays the identifying equality of hom groups and reads K-injectivity off it.\par\medskip
\noindent\textit{Editorial note.} Cross-chapter restatement: the same label occurs in chapter cohomology and in chapter more-algebra. Those statements are the special cases of this one, so per the lane-s4 rule only this more general occurrence is counted here.\par\medskip
\noindent\textit{Editorial note.} The single labelled reference occurring inside the counted spans points to a chapter of the Project not staged in this card's source tree; it is carried in assembly.external\_references as a non-load-bearing navigation pointer recorded at the invoking line.\par\medskip
\noindent\textit{Editorial note.} The preamble is the chapter's own retained minimal preamble (sources/stacks-sites-cohomology/retained-preamble.tex), a verbatim excerpt of the native preamble: the Xy-pic setup and the native theorem declarations with the mathematical macros.\par\medskip
\noindent\textit{Allowed standard prerequisites (editorial).}\begin{itemize}
\item K-injective complex of O-modules. Same chapter, the defining property of K-injectivity in terms of the vanishing of homotopy classes of maps into it from acyclic complexes. Supplies the property to be verified for the restricted internal hom over O'.
\item Restriction of scalars along a map of sheaves of rings. Same chapter, the functor SheafHom\_O(O', -) sending O-modules to O'-modules and the induced map of sites. Supplies the complex whose K-injectivity is asserted.
\item Adjointness of the restriction functor. Modules on Sites, the identification Hom\_\{K(O')\}(G\textasciicircum{}bullet, Hom\_O(O', I\textasciicircum{}bullet)) = Hom\_\{K(O)\}(G\textasciicircum{}bullet, I\textasciicircum{}bullet). Supplies the equality of homotopy groups from which K-injectivity is read off.
\end{itemize}\medskip
\ifdefined\hypertarget\hypertarget{sourcepointer-1}{}\fi
\noindent\textit{Source reference.} Original lemma 0932: sources/stacks-sites-cohomology/pinned-original-source.tex, source lines 4629--4629. Cross-chapter navigation pointer invoked at proof lines 4628-4629, where the author uses the adjunction identifying homomorphisms into the restricted internal hom with homomorphisms over the smaller sheaf of rings. The target chapter is not staged in this card's source tree, so the pointer records the invoking line; it only resolves the reference for the reader, and no bibliography citation is involved.\par\medskip
\begin{lemma}
\label{lemma-hom-K-injective}
Let $\mathcal{C}$ be a site. Let $\mathcal{O} \to \mathcal{O}'$ be a map
of sheaves of rings. If $\mathcal{I}^\bullet$ is a K-injective complex of
$\mathcal{O}$-modules, then
$\SheafHom_\mathcal{O}(\mathcal{O}', \mathcal{I}^\bullet)$
is a K-injective complex of $\mathcal{O}'$-modules.
\end{lemma}

\begin{proof}
This is true because
$\Hom_{K(\mathcal{O}')}(\mathcal{G}^\bullet,
\Hom_\mathcal{O}(\mathcal{O}', \mathcal{I}^\bullet)) =
\Hom_{K(\mathcal{O})}(\mathcal{G}^\bullet, \mathcal{I}^\bullet)$
by Modules on Sites, Lemma \ref{sites-modules-lemma-adjoint-hom-restrict}.
\end{proof}

\end{document}
